\subsection{\texorpdfstring{\(\mathbf{R}^n\) 中的级数}{R 的 n 次幂中的级数}}

通项为 \(\boldsymbol{x}_{m} = (x_{mi})_{1 \leqslant i \leqslant n}\) 的级数在 \(R^{n}\) 中收敛，当且仅当 n 个级数 \((x_{mi})_{m \in \mathbb{N}}\) 中的每一个都在 R 中收敛。

定义 1. 若 \(\mathbf{R}^n\) 中点的一个级数的各项的欧氏范数所成的级数收敛，则称该级数是绝对收敛的。

命题 4. \(\mathbf{R}^n\) 中点的一个级数是交换收敛的，当且仅当它是绝对收敛的。

这是由第 III 章 § 5 第 7 小节命题 9 和上面的定理 1 得出的。

第 IV 章 § 7 中给出的例子表明， \(\mathbf{R}^n\) 中的一个级数可以收敛而非绝对收敛。
% semantic-fragments: legacy-input-seam
\relax{}
\ExerciseGroupsStart
